\section*{Problemas}

\subsection*{Enunciados de los problemas}

% Recordar
\begin{problema}
	\label{prob:ca86634}
	Para validar la precisión del pronóstico de demanda de electricidad residencial en una subestación urbana ($Y$, en megavatios-hora), un ingeniero somete a escrutinio un modelo de regresión sobre una muestra de validación de $n_{\text{test}} = 4$ días independientes:
	\begin{center}
		\begin{tabular}{ccc}
			\hline
			\textbf{Día ($i$)} & \textbf{Consumo real ($Y_i$)} & \textbf{Predicción del modelo ($\hat{Y}_i$)} \\ \hline
			1 & 120.0 & 115.0 \\
			2 & 135.0 & 138.0 \\
			3 & 110.0 & 112.0 \\
			4 & 145.0 & 141.0 \\ \hline
		\end{tabular}
	\end{center}
	\begin{enumerate}
		\item Calcule el \textbf{Error Cuadrático Medio de Prueba} ($\text{MSE}_{\text{test}} = \frac{1}{n_{\text{test}}}\sum(Y_i-\hat{Y}_i)^2$).
		\item Calcule la \textbf{Raíz del Error Cuadrático Medio} ($\text{RMSE}_{\text{test}}$) y el \textbf{Error Absoluto Medio} ($\text{MAE}_{\text{test}}$).
		\item Explique por qué $\text{RMSE}_{\text{test}}$ resulta invariablemente mayor o igual que $\text{MAE}_{\text{test}}$ y qué penalización geométrica impone el primero ante errores de pronóstico grandes.
	\end{enumerate}
\end{problema}

% Comprender
\begin{problema}
	\label{prob:7d813ee}
	Un equipo de ciencia de datos ajusta un modelo de regresión múltiple para estimar el valor comercial de inmuebles residenciales ($Y$, en miles de dólares). Con $n = 200$ observaciones ($80\%$ entrenamiento, $20\%$ prueba): $R_{\text{train}}^2 = 0.760$, $RSE_{\text{train}} = 24.50$; $R_{\text{test}}^2 = 0.720$, $RSE_{\text{test}} = 27.80$.
	\begin{enumerate}
		\item Evalúe si existe evidencia de sobreajuste comparando la brecha entre las métricas intra-muestra y extra-muestra.
		\item Si el valor comercial promedio en prueba es $\bar{Y}_{\text{test}} = 250.0$ miles, calcule el porcentaje medio de error predictivo residual respecto a la media del mercado.
		\item Argumente por qué la división simple en dos subconjuntos (\emph{Holdout Method}) puede ser sensible o inestable con muestras moderadas, y proponga una alternativa técnica superadora.
	\end{enumerate}
\end{problema}

% Aplicar
\begin{problema}
	\label{prob:d0d57bc}
	Al auditar un modelo econométrico lineal con $p = 15$ variables sobre $n = 120$ transacciones corporativas, separadas en $80\%$ calibración ($n_{\text{train}}=96$) y $20\%$ evaluación ($n_{\text{test}}=24$):
	\begin{center}
		\begin{tabular}{|l|c|c|}
			\hline
			\textbf{Métrica} & \textbf{Calibración} & \textbf{Evaluación} \\ \hline
			$R^2$ & 0.850 & 0.620 \\
			$RSE$ & 15.20 & 28.70 \\
			$MAE$ & 12.10 & 22.50 \\ \hline
		\end{tabular}
	\end{center}
	\begin{enumerate}
		\item Calcule las razones de degradación predictiva extra-muestra $RSE_{\text{test}}/RSE_{\text{train}}$ y $R_{\text{test}}^2/R_{\text{train}}^2$.
		\item Realice el diagnóstico diferencial: ¿presenta el modelo \textbf{sobreajuste extremo (alta varianza)} o \textbf{subajuste (alto sesgo)}? Justifique con los umbrales estándar de la industria.
		\item Elabore un plan de rectificación con al menos cuatro intervenciones técnicas precisas (regularización, selección de variables, validación cruzada anidada, etc.).
	\end{enumerate}
\end{problema}

% Analizar
\begin{problema}
	\label{prob:cb39746}
	Para modelar la elasticidad precio-demanda del transporte público ($n = 100$ trayectos), se calibraron tres modelos con distinta complejidad ($p$ incluye intercepto y $\sigma^2$):
	\begin{center}
		\begin{tabular}{|l|c|c|c|}
			\hline
			\textbf{Modelo} & \textbf{$p$} & \textbf{SCE} & \textbf{SCE$/n$} \\ \hline
			A (Básico) & 3 & 450.00 & 4.500 \\
			B (Intermedio) & 5 & 380.00 & 3.800 \\
			C (Saturado) & 8 & 350.00 & 3.500 \\ \hline
		\end{tabular}
	\end{center}
	Use $\text{AIC} = n\ln(\text{SCE}/n)+2p$, $\text{BIC} = n\ln(\text{SCE}/n)+p\ln(n)$, con $\ln(4.50)=1.5041$, $\ln(3.80)=1.3350$, $\ln(3.50)=1.2528$, $\ln(100)=4.6052$.
	\begin{enumerate}
		\item Calcule manualmente los valores de AIC y BIC para los tres modelos.
		\item Determine qué modelo resulta seleccionado bajo el criterio AIC y cuál bajo BIC.
		\item Explique por qué el BIC aplica una penalización asintóticamente más severa sobre la complejidad ($p$) que el AIC cuando $n\ge8$, y discuta las implicaciones para la selección de modelos en \emph{Big Data}.
	\end{enumerate}
\end{problema}

% Evaluar
\begin{problema}
	\label{prob:2b2e597}
	Para optimizar un sistema predictivo de ventas minoristas, se comparan dos arquitecturas bajo Validación Cruzada de 10 pliegues:
	\begin{itemize}
		\item \textbf{Modelo A ($p_A=3$):} $\bar{R}_{\text{CV},A}^2 = 0.780 \pm 0.040$ (media $\pm$ desviación estándar entre pliegues).
		\item \textbf{Modelo B ($p_B=5$):} $\bar{R}_{\text{CV},B}^2 = 0.800 \pm 0.060$.
	\end{itemize}
	\begin{enumerate}
		\item Determine qué modelo exhibe mayor poder explicativo medio y cuál demuestra mayor estabilidad mediante el Coeficiente de Variación porcentual.
		\item Al desplegar ambos en producción sobre $n=150$ tiendas, calcule el $R_{\text{adj}}^2$ esperado de cada uno.
		\item Evalúe, aplicando el Principio de Parsimonia (Navaja de Ockham), qué modelo debe adoptar la gerencia, sopesando la ganancia marginal en ajuste contra el costo de complejidad y la mayor inestabilidad del modelo saturado.
	\end{enumerate}
\end{problema}

% Crear
\begin{problema}
	\label{prob:9689172}
	Diseñe un problema original de validación de modelos (distinto de inmuebles, electricidad, transacciones corporativas, transporte público o ventas minoristas ya vistos): un equipo de ciencia de datos en una plataforma de streaming evalúa un modelo de predicción de cancelación de suscripción (\emph{churn}). Divide $n=500$ clientes en $350$ de entrenamiento y $150$ de prueba, obteniendo $R_{\text{train}}^2=0.82$, $RSE_{\text{train}}=8.5$; $R_{\text{test}}^2=0.79$, $RSE_{\text{test}}=9.6$. Calcule la brecha $\Delta R^2$ y el incremento porcentual relativo del $RSE$, y evalúe si el modelo exhibe sobreajuste leve/admisible o severo/patológico según la escala estándar ($\Delta R^2\le0.05$ e incremento de $RSE\le15\%$ es admisible; $\Delta R^2>0.15$ o incremento $>40\%$ es severo).
\end{problema}

\subsection*{Soluciones detalladas}

% Recordar
\begin{solproblema}[prob:ca86634]
	\begin{enumerate}
		\item Los residuos son $e_1=+5.0$, $e_2=-3.0$, $e_3=-2.0$, $e_4=+4.0$, con $\sum e_i^2=25+9+4+16=54.00$. Entonces $\text{MSE}_{\text{test}}=54.00/4=13.50$ (MWh)$^2$.
		\item $\text{RMSE}_{\text{test}}=\sqrt{13.50}\approx3.6742$ MWh. Con $\sum|e_i|=5+3+2+4=14.00$: $\text{MAE}_{\text{test}}=14.00/4=3.5000$ MWh.
		\item $\text{RMSE}\ge\text{MAE}$ siempre, por la desigualdad entre la norma cuadrática media ($L_2$) y la norma absoluta media ($L_1$), con igualdad solo si todos los $|e_i|$ son idénticos. Al elevar al cuadrado cada residuo antes de promediar, el RMSE penaliza desproporcionadamente los errores grandes: un error de $10$ MWh contribuye $100$ unidades a la suma, equivalente a cien errores individuales de $1$ MWh. Por eso el RMSE se prefiere cuando los errores grandes son especialmente costosos.
	\end{enumerate}
\end{solproblema}

% Comprender
\begin{solproblema}[prob:7d813ee]
	\begin{enumerate}
		\item $\Delta R^2 = 0.760-0.720=0.040$ (4.0 puntos). Incremento relativo de RSE: $(27.80-24.50)/24.50\times100\approx13.47\%$. Con la escala estándar ($\Delta R^2\le0.05$ e incremento de RSE $\le15\%$ es leve/admisible), \textbf{no existe sobreajuste patológico}: el modelo generaliza razonablemente.
		\item Error de pronóstico relativo $=RSE_{\text{test}}/\bar{Y}_{\text{test}}\times100=27.80/250.00\times100=11.12\%$.
		\item Con $n=200$ moderado, reservar un único bloque fijo de $n_{\text{test}}=40$ genera alta varianza (la estimación depende de qué observaciones cayeron por azar en el bloque de prueba) y pérdida de potencia muestral (los parámetros se calibran con menos datos). La alternativa superadora es la \textbf{Validación Cruzada de $k$ pliegues} ($k=5$ o $k=10$), que rota las particiones para que cada dato se use exactamente una vez como prueba, promediando y estabilizando la estimación del error de generalización.
	\end{enumerate}
\end{solproblema}

% Aplicar
\begin{solproblema}[prob:d0d57bc]
	\begin{enumerate}
		\item $RSE_{\text{test}}/RSE_{\text{train}} = 28.70/15.20 \approx1.8882$ (el error casi se duplica, $+88.82\%$). $R_{\text{test}}^2/R_{\text{train}}^2 = 0.620/0.850\approx0.7294$ (caída del $27.06\%$ en varianza explicada).
		\item El modelo exhibe \textbf{Sobreajuste Extremo (Alta Varianza)}: coexiste un desempeño artificialmente alto en entrenamiento ($R_{\text{train}}^2=0.850$) con un colapso extra-muestra (ratio $RSE>1.5$, pérdida de $R^2>0.20$). Con $p=15$ predictores sobre solo $n_{\text{train}}=96$ (apenas $6.4$ observaciones por variable, muy por debajo del umbral de seguridad de $15$–$20$), el modelo memorizó ruido idiosincrásico en vez de patrones generalizables.
		\item Plan de rectificación: (1) regularización Ridge/Lasso para contraer los coeficientes; (2) selección de variables (eliminación recursiva o VIF$>5$) para reducir a 4–6 predictores de alta señal; (3) validación cruzada anidada (bucle interno para hiperparámetros, bucle externo para evaluación) para evitar fuga de información; (4) transformaciones (log, Box-Cox) y regresión robusta para amortiguar valores atípicos de alto apalancamiento.
	\end{enumerate}
\end{solproblema}

% Analizar
\begin{solproblema}[prob:cb39746]
	\begin{enumerate}
		\item Modelo A: $\text{AIC}_A=100(1.5041)+2(3)=150.41+6.00=156.41$; $\text{BIC}_A=150.41+3(4.6052)=150.41+13.82=164.23$.
		Modelo B: $\text{AIC}_B=100(1.3350)+2(5)=133.50+10.00=143.50$; $\text{BIC}_B=133.50+5(4.6052)=133.50+23.03=156.53$.
		Modelo C: $\text{AIC}_C=100(1.2528)+2(8)=125.28+16.00=141.28$; $\text{BIC}_C=125.28+8(4.6052)=125.28+36.84=162.12$.
		\item El criterio de selección es el menor valor numérico. Bajo \textbf{AIC}: el \textbf{Modelo C} ($141.28$) es óptimo. Bajo \textbf{BIC}: el \textbf{Modelo B} ($156.53$) es óptimo.
		\item El AIC penaliza cada parámetro adicional con una constante fija $+2$, mientras que el BIC penaliza con $+\ln(n)$, que supera a $2$ tan pronto $n\ge8$ ($\ln(8)\approx2.079$). Para $n=100$, cada parámetro extra cuesta $+4.61$ en BIC frente a $+2.00$ en AIC —más del doble—, por lo que el BIC descarta al Modelo C (saturado) mientras el AIC lo prefiere. En \emph{Big Data} ($n$ muy grande, $\ln n$ enorme), el AIC tiende crónicamente a sobreseleccionar variables con efectos ínfimos, mientras que el BIC posee la propiedad de consistencia asintótica: converge en probabilidad al modelo verdadero cuando $n\to\infty$, siendo preferible para selección de modelos parsimoniosos a gran escala.
	\end{enumerate}
\end{solproblema}

% Evaluar
\begin{solproblema}[prob:2b2e597]
	\begin{enumerate}
		\item El Modelo B tiene mayor poder explicativo bruto ($0.800$ vs. $0.780$). Coeficientes de variación: $CV_A=0.040/0.780\times100\approx5.13\%$, $CV_B=0.060/0.800\times100\approx7.50\%$. El \textbf{Modelo A es más estable} (menor CV inter-pliegues).
		\item Con $n=150$: $R_{\text{adj,A}}^2 = 1-\left[\frac{149}{146}(0.220)\right]\approx0.7755$; $R_{\text{adj,B}}^2 = 1-\left[\frac{149}{144}(0.200)\right]\approx0.7931$.
		\item La ganancia neta ajustada del Modelo B sobre A es de apenas $+1.76$ puntos porcentuales, a cambio de $66.7\%$ más variables y una inestabilidad inter-pliegues considerablemente mayor ($CV$ un $46\%$ más alto). Bajo la Navaja de Ockham —preferir la hipótesis más simple entre las que explican los datos con similar exactitud—, \textbf{se recomienda el Modelo A}: su simplicidad, estabilidad y precisión virtualmente idéntica protegen a la corporación del sobreajuste sin sacrificar poder predictivo relevante.
	\end{enumerate}
\end{solproblema}

% Crear
\begin{solproblema}[prob:9689172]
	\begin{align}
		\Delta R^2 = R_{\text{train}}^2 - R_{\text{test}}^2 = 0.82-0.79 = 0.03 \text{ (3.0 puntos porcentuales).}
	\end{align}
	\begin{align}
		\text{Incremento relativo de } RSE = \frac{9.6-8.5}{8.5}\times100 = \frac{1.1}{8.5}\times100 \approx 12.94\%.
	\end{align}
	Como $\Delta R^2 = 0.03 \le 0.05$ y el incremento relativo del RSE ($12.94\%$) es menor al umbral de $15\%$, el modelo se ubica en la categoría de \textbf{sobreajuste leve / admisible}, propio de la fluctuación muestral normal en una muestra de prueba de $n_{\text{test}}=150$ clientes. El modelo de predicción de \emph{churn} generaliza razonablemente bien a clientes nuevos y puede desplegarse en producción sin requerir intervenciones de regularización adicionales, aunque conviene monitorear su desempeño periódicamente conforme cambie el comportamiento de los suscriptores.
\end{solproblema}
